\documentclass[11pt]{article}
\usepackage[margin=0.82in]{geometry}
\usepackage[T1]{fontenc}
\usepackage{lmodern,amsmath,amssymb,amsthm,booktabs}
\usepackage[colorlinks=true,linkcolor=blue,urlcolor=blue,citecolor=blue]{hyperref}
\hypersetup{pdftitle={A universal 3-adic limit for Fibonacci-Collatz sequences},pdfauthor={Blake Stone},pdfsubject={Proofs of the OEIS A105801 and A181717 residue conjectures}}
\setlength{\parindent}{0pt}
\setlength{\parskip}{0.45em}
\newtheorem{theorem}{Theorem}
\newtheorem{lemma}[theorem]{Lemma}
\newtheorem{corollary}[theorem]{Corollary}
\newcommand{\N}{\mathbb N}
\newcommand{\Z}{\mathbb Z}
\title{A universal $3$-adic limit for Fibonacci--Collatz sequences\\
\large Resolving the residue conjectures in OEIS A105801 and A181717}
\author{Blake Stone\\\normalsize Proof developed with OpenAI Astra 6.2}
\date{November 26, 2026}
\begin{document}
\maketitle
\begin{abstract}
For the recurrence that applies a Collatz step to the sum of the two previous
terms, we prove that every pair of distinct nonnegative starting values gives
a sequence converging $3$-adically to $-1/5$. Equivalently, for every $k\ge0$,
all sufficiently late terms satisfy $3^k\mid 5a_n+1$. This proves Giovanni
Resta's residue-stabilization conjectures for A105801 and A181717 and gives
the stable residue explicitly. The proof uses an affine change of variables,
integer descent, and repeated increases of common divisibility by $3$.
\end{abstract}

\section{Statement and change of variables}

For nonnegative integers $x,y$, define $a_1=x$, $a_2=y$, and, for $n\ge1$,
\begin{equation}
 a_{n+2}=F(a_n,a_{n+1}),\qquad
 F(u,v)=
 \begin{cases}
 (u+v)/2,&u+v\text{ even},\\
 3(u+v)+1,&u+v\text{ odd}.
 \end{cases}
 \label{eq:recurrence}
\end{equation}
We call these the averaging and expanding branches, respectively.
The seeds $(1,2)$ give OEIS A105801~\cite{A105801}; the seeds $(0,1)$ give
A181717~\cite{A181717}. Both entries record Resta's November 17, 2010
conjecture that the sequence eventually has a constant residue modulo
every power of $3$.

\begin{theorem}\label{thm:main}
If $x,y\in\N$ and $x\ne y$, then for every integer $k\ge0$ there exists
$N_k$ such that
\begin{equation}
 3^k\mid 5a_n+1\qquad\text{for every }n\ge N_k.
 \label{eq:divisibility}
\end{equation}
Consequently, every such sequence eventually has the same residue
$-5^{-1}\pmod{3^k}$, independently of its starting values.
\end{theorem}

The decisive change of variables is
\begin{equation}
 b_n=5a_n+1.
 \label{eq:change}
\end{equation}
The branch condition is unchanged because $b_n+b_{n+1}=5(a_n+a_{n+1})+2$
has the same parity as $a_n+a_{n+1}$. Direct substitution gives
\begin{equation}
 b_{n+2}=
 \begin{cases}
 (b_n+b_{n+1})/2,&\text{averaging branch},\\
 3(b_n+b_{n+1}),&\text{expanding branch}.
 \end{cases}
 \label{eq:transformed}
\end{equation}
Every $b_n$ is a positive integer. The additive constant in
\eqref{eq:recurrence} has disappeared.

\section{Parity and the recurrence of two expanding steps}

\begin{lemma}\label{lem:distinct}
Consecutive terms remain unequal. Moreover, no tail of the sequence can
consist entirely of averaging steps.
\end{lemma}
\begin{proof}
If $u\ne v$ and the averaging branch applies, then $F(u,v)=v$ would imply
$u=v$. On the expanding branch, $F(u,v)=3u+3v+1>v$. Thus inequality of
consecutive terms is preserved by induction.

Set $d_n=|a_{n+1}-a_n|$. These are positive integers. At every averaging
step,
\[
 d_{n+1}
 =\left|\frac{a_n+a_{n+1}}2-a_{n+1}\right|
 =\frac{d_n}{2}.
\]
An infinite averaging tail would therefore give an infinite strictly
decreasing sequence of positive integers, which is impossible.
\end{proof}

\begin{lemma}\label{lem:blocks}
There are arbitrarily late indices $j$ for which $a_{j-1}$ is even and
$a_j$ is odd. Each such index forces two consecutive expanding steps.
\end{lemma}
\begin{proof}
By Lemma~\ref{lem:distinct}, some expanding step occurs. Its output is
even. From that point onward, each consecutive pair contains an even term:
an expanding step outputs an even term, while an averaging step from a pair
containing an even term must start from two even terms, leaving the previous
even term in the next pair.

Odd terms occur arbitrarily late. Otherwise an all-even tail would have only
averaging steps, contradicting Lemma~\ref{lem:distinct}. Every sufficiently
late odd term $a_j$ has an even predecessor. Hence $a_{j-1}+a_j$ is odd,
so $a_{j+1}$ is produced by an expanding step and is even. Then
$a_j+a_{j+1}$ is odd as well, forcing a second expanding step.
\end{proof}

This argument also handles two odd starting values: the even-term invariant
is established after the first expanding step, rather than assumed at the
initial pair.

\section{Divisibility increases without bound}

\begin{lemma}\label{lem:preserve}
If $3^k$ divides two consecutive transformed terms $b_n,b_{n+1}$, then it
divides every subsequent transformed term. If the next two steps are
expanding, both $b_{n+2}$ and $b_{n+3}$ are divisible by $3^{k+1}$.
\end{lemma}
\begin{proof}
For an averaging step, $2b_{n+2}=b_n+b_{n+1}$ is divisible by $3^k$.
Since $2$ is coprime to $3^k$, the same is true of $b_{n+2}$. For an
expanding step, $b_{n+2}=3(b_n+b_{n+1})$ is divisible by $3^{k+1}$ and
therefore by $3^k$. This proves persistence by induction along the sequence.

If two expanding steps follow, then
\[
 b_{n+2}=3(b_n+b_{n+1}),\qquad
 b_{n+3}=3(b_{n+1}+b_{n+2}).
\]
Both right-hand sides are divisible by $3^{k+1}$.
\end{proof}

\begin{proof}[Proof of Theorem~\ref{thm:main}]
We induct on $k$. The case $k=0$ is immediate. Suppose all terms $b_n$
from index $N_k$ onward are divisible by $3^k$. By
Lemma~\ref{lem:blocks}, choose an odd term $a_j$ with even predecessor and
$j-1\ge N_k$. The next two steps are expanding. Lemma~\ref{lem:preserve}
therefore gives $3^{k+1}\mid b_{j+1},b_{j+2}$, and its persistence statement
extends this to every later term. We may take $N_{k+1}=j+1$.

Finally, $5$ is invertible modulo $3^k$. Thus $3^k\mid 5a_n+1$ determines
a unique residue for $a_n$, namely $-5^{-1}\pmod{3^k}$.
\end{proof}

\section{The stable residue and the \texorpdfstring{$3$}{3}-adic interpretation}

An explicit representative is
\begin{equation}
 c_k=2^k\bmod5,\qquad
 r_k=\frac{c_k3^k-1}{5}
     =\left\lfloor\frac{c_k3^k}{5}\right\rfloor.
 \label{eq:residue}
\end{equation}
Indeed $1\le c_k\le4$ and $c_k3^k\equiv6^k\equiv1\pmod5$, so $r_k$ is
an integer satisfying $0\le r_k<3^k$ and $5r_k+1=c_k3^k$.
The coefficients $c_k$ have period $4$:
\[
 c_k=1,2,4,3\quad\text{as }k\equiv0,1,2,3\pmod4.
\]

The following examples use OEIS indexing. Each displayed $N$ is the first
index at which two adjacent terms are congruent to $r_k$; the persistence
lemma then certifies every term from $N$ onward. The finite calculations are
reproduced by \texttt{examples.py} and are illustrations of the theorem.
\begin{center}
\begin{tabular}{rrrrr}
\toprule
$k$ & $3^k$ & $r_k$ & $N$ for A105801 & $N$ for A181717\\
\midrule
1&3&1&6&2\\
2&9&7&10&7\\
3&27&16&13&10\\
4&81&16&15&15\\
5&243&97&19&19\\
6&729&583&22&23\\
7&2187&1312&24&25\\
8&6561&1312&28&30\\
\bottomrule
\end{tabular}
\end{center}

In the $3$-adic metric, \eqref{eq:divisibility} says that
$|a_n+1/5|_3\le3^{-k}$ eventually for each $k$, since $|5|_3=1$.
Thus
\[
 \boxed{\displaystyle a_n\longrightarrow-\frac15\quad\text{in }\Z_3.}
\]
The compatible base-$3$ digits repeat the block $0121$:
$-1/5=16\sum_{j\ge0}81^j$ in $\Z_3$.
If the two starting values are equal to $x$, the sequence instead remains
constant at $x$. This accounts for every pair of nonnegative starting values.

\section{Formal verification and relation to earlier work}

The accompanying \texttt{Proof.lean} formalizes the original recurrence,
descent, parity arguments, divisibility persistence and lifting, and the
eventual residue statement for every distinct pair of nonnegative seeds.
It also proves \eqref{eq:residue} and specializes the general theorem to both
OEIS sequences as \texttt{A105801.conjecture} and \texttt{A181717.conjecture}.
The complete package passed a fresh Lean 4.19.0 build. It imports only
\texttt{Std}, has no external dependencies, and uses only the standard
logical axioms \texttt{propext} and \texttt{Quot.sound}. There are no
unfinished proofs, added axioms, or \texttt{native\_decide} calls.
Build commands, axiom reports, and a separate review are supplied.

The $3$-adic formulation above is the standard interpretation of the proved
arbitrary-power congruences. The formal result is stated directly using
natural-number divisibility and residues, without requiring $3$-adic
topological machinery.

Greene~\cite{Greene} studied a larger family of integer difference equations.
The transformed recurrence \eqref{eq:transformed} is its case
$c_1=c_2=3$, $d=6$; his result already gives ordinary divergence for unequal
positive starting values. That result concerns size, while the theorem here
establishes eventual stabilization modulo every $3^k$ and identifies the
common $3$-adic limit. A 2025 Math.StackExchange discussion~\cite{MSE}
also observes the persistence of the residue $7$ modulo $9$.

The retrieved OEIS entries still state the arbitrary-power claims as
conjectures. Searches by identifier and equivalent recurrence found no
earlier proof of those claims. This records the scope of the source check;
it does not establish exhaustive publication priority or OEIS acceptance.

\begin{thebibliography}{9}
\bibitem{A105801} Z. Seidov, OEIS A105801, September 12, 2006;
conjecture by G. Resta, November 17, 2010.
\url{https://oeis.org/A105801}.
\bibitem{A181717} R. Stephan, OEIS A181717, November 17, 2010;
conjecture by G. Resta, same date.
\url{https://oeis.org/A181717}.
\bibitem{Greene} J. Greene, ``The Unboundedness of a Family of Difference
Equations Over the Integers,'' \emph{The Fibonacci Quarterly} 46/47(2)
(2008/2009), 146--152.
\url{https://www.fq.math.ca/Papers1/46_47-2/Greene_1-09.pdf}.
\bibitem{MSE} Exit load, question on Fibonacci--Collatz digital roots,
with a comment by Blue, January 27--28, 2025.
\url{https://math.stackexchange.com/questions/5028237/}.
\end{thebibliography}
\end{document}
